% New manuscript insert: a derived implication from the existing G18 inputs.
% Not automatically input into the main manuscript. No upstream status is changed.
% Requires amsmath, amssymb, and theorem/lemma/proof environments.
\subsection{Uniform relative carrier isolation on the punctured zone}
\label{sec:g18-relative-gap-bridge}

We retain the fixed-spacing physical-band identification and the
coupling-analytic weighted symbol of Lemma~\ref{lem:g18-u-analytic-symbol}.
The conclusion below strengthens patchwise isolation; it is not an
independent proof of those inputs. Write
\[
 L(k)=B(k)B(k)^\dagger=q(k)I-w(k)w(k)^\dagger,\qquad
 q(k)=4\sum_{j=1}^3\sin^2(k_j/2),\qquad t_3=\frac5{612}.
\]
The symmetry-covariant frame has the irreducible cubic $T_1$ representation
at $\Gamma$, where inversion is scalar. Thus
$h_u(0)=E_\Gamma(u)I$ and $\nabla_kh_u(0)=0$.
The latter assertion does not require entrywise evenness in a
plaquette-base convention. If
$h_u(-k)=J(k)h_u(k)J(k)^{-1}$ with $J(0)=pI$, differentiating at zero
leaves the commutator with $E_\Gamma I$, which vanishes, and gives
$-\partial_jh_u(0)=\partial_jh_u(0)$.
Gram orthonormalization preserves covariance by functional calculus.

\begin{lemma}[Weighted centering]
\label{lem:g18-weighted-centering}
Let $R\in\mathcal A_\mu$ and assume $\nabla R(0)=0$. Then, on
$[-\pi,\pi]^3$,
\[
 \|R(k)-R(0)\|\le\gamma_\mu\|R\|_{\mu,\sharp}q(k),\qquad
 \gamma_\mu=\frac{\pi^2}{2e^2\mu^2}.
\]
\end{lemma}
\begin{proof}
Exponential summability permits differentiation of the Fourier series.
The zero first moment gives
\[
 R(k)-R(0)=\sum_xR(x)(e^{ik\cdot x}-1-ik\cdot x).
\]
The integral Taylor remainder satisfies
$|e^{it}-1-it|\le t^2/2$. Apply it entrywise, use
$|k\cdot x|\le |k|_2|x|_1$, and apply the row/column Schur bound:
\[
 \|R(k)-R(0)\|\le\frac{|k|_2^2}{2}
 \sup_{s\ge0}(s^2e^{-\mu s})\|R\|_{\mu,\sharp}.
\]
The supremum is $4/(e^2\mu^2)$ and
$|k|_2^2\le(\pi^2/4)q(k)$, proving the assertion.
\end{proof}

\begin{theorem}[One coupling domain for every nonzero momentum]
\label{thm:g18-global-relative-gap}
Assume that $u\mapsto h_u\in\mathcal A_\mu$ is holomorphic on
$|u|<\rho_0$, Hermitian for real $u$, and has the symmetry properties above.
Assume also the full matrix coefficient identity
\[
 h_u(k)=s_{\le2}(u)I+t_3u^2L(k)+O_{\mathcal A_\mu}(u^3).
\]
Choose $0<\rho<\rho_0$ within all upstream smallness domains and set
\[
 M_\rho=\sup_{|\zeta|=\rho}\|h_\zeta\|_{\mu,\sharp},\qquad
 K=\frac{\pi^2M_\rho}{e^2\mu^2\rho^3},\qquad
 u_* = \min\left\{\rho/2,\frac{t_3}{4K}\right\}.
\]
If $K=0$, omit its restriction. For every real $0<|u|<u_*$ and every
$k\ne0$, the lowest eigenvalue of $h_u(k)$ is simple, and
\begin{align}
 \|h_u(k)-E_\Gamma(u)I-t_3u^2L(k)\|&\le K|u|^3q(k),
 \label{eq:g18-relative-tail}\\
 E_1(k,u)-E_0(k,u)&\ge u^2q(k)(t_3-2K|u|)
 \ge\tfrac12t_3u^2q(k).
 \label{eq:g18-relative-gap-global}
\end{align}
The coupling domain is independent of momentum.
\end{theorem}
\begin{proof}
Cauchy's estimate in $\mathcal A_\mu$ bounds the Taylor tail beginning at
order three by $2M_\rho|u|^3/\rho^3$ for $|u|\le\rho/2$.
Its first momentum derivative vanishes, coefficientwise by symmetry and
holomorphy. Lemma~\ref{lem:g18-weighted-centering}, with its exact Gamma
value subtracted, gives~\eqref{eq:g18-relative-tail}.
The comparison eigenvalues after subtracting $E_\Gamma I$ are
$0,t_3u^2q,t_3u^2q$. Weyl's eigenvalue inequality shifts each by at most
$K|u|^3q$. Subtracting the two worst shifts proves
\eqref{eq:g18-relative-gap-global}.
\end{proof}

Let $P_0=ww^\dagger/q$ and let $P_u$ denote the exact lowest projection.
A useful companion estimate is
\[
 \|P_u(k)-P_0(k)\|\le\frac{K|u|}{t_3-2K|u|},\qquad k\ne0.
\]
Indeed, put $a=t_3u^2q$ and $\varepsilon=K|u|^3q$.
The exact eigenvalue relative to the Gamma scalar lies in
$[-\varepsilon,\varepsilon]$. Projecting its eigenvector equation onto
$1-P_0$ gives
$\|(1-P_0)v\|\le\varepsilon\|P_0v\|/(a-2\varepsilon)$,
and the rank-one projection distance is $\|(1-P_0)v\|$.
After shrinking the common interval, the distance is less than $1/2$ and
\[
 v_u(k)=\frac{P_u(k)w(k)/\sqrt{q(k)}}
 {\|P_u(k)w(k)/\sqrt{q(k)}\|}
\]
is a global periodic real-analytic unit section on the punctured zone.
Its direct integral is a reducing carrier subspace. This is a fiberwise
energy--momentum construction; no energy-only isolating contour is asserted.

For the normalized three-orientation Wilson synthesis $S_W$, the
complete-band frame bound $S_W^\dagger S_W\succeq z_*I$ implies
\[
 \sum_{\alpha=1}^3|\langle v_u,S_We_\alpha\rangle|^2\ge z_*.
\]
The left and right Gram matrices of the square synthesis map have the same
least eigenvalue. The usual normalized cube-source estimate also becomes
uniform:
\[
 \|P_uS_Wv_0\|\ge1-\|P_u-P_0\|-\|S_W-I\|,\qquad v_0=w/\sqrt q.
\]
Neither statement makes the compact cube boundary a uniformly normalized
local source: its coefficient is $w$, not $w/\sqrt q$, and its weight
vanishes with $q$ near Gamma.

\paragraph{Boundary.}
The exact absolute internal gap still tends to zero at Gamma; no
rank-one choice there is claimed. The constants are existential and the
finite-order threshold $\sqrt{2/51}$ is not an all-orders threshold.
The theorem inherits the G18 analytic-symbol and physical-band inputs.
It proves no finite-temporal-spacing Wilson identification or continuum
limit. The patch theorem remains valid, but deterioration of its coupling
threshold with the patch minimum of $q$ is not necessary under the stronger
weighted hypotheses used here.
